\section{The precise contextad connection}\label{sec:contextads}
Contextads organize the operational side of this comparison, provided their
actual structural hypotheses are supplied. They do not extract an
observation doctrine from a table. In the current source, Capucci--Myers
\cite[Definition 5.25, Examples 5.33--5.35]{CM2}, a contextad has indexed
grades $K:C^{\mathrm{op}}\to\Cat$, extension $e:\int K\to C$,
unit grades and coherent accumulation. A contextual arrow $(P,f):A\Rightarrow B$
has body $f:e_AP\to B$. The grade of a composite with $(R,g):B\Rightarrow D$
is $P\otimes f^*R$.

A strong contextad morphism supplies $F:C\to C'$, a cartesian functor
$G:\int K\to\int K'$ over $F$, and a natural lineator
\[
 \ell_{A,P}:e'_{FA}(G_AP)\longrightarrow F(e_AP),
\]
together with coherent invertible unit and multiplication comparisons
\[
 G_AI_A\cong I'_{FA},\qquad
 G_A(P\otimes Q)\cong
 G_AP\otimes'\ell_{A,P}^*G_{e_AP}Q.
\]
The cartesian condition means preservation of cartesian lifts and hence
$G(u^*R)\cong(Fu)^*G(R)$; it does not say that the entire square of
categories is a pullback. Compatibility with the extension structure,
units and associativity is required. The lineator itself need not be
invertible. The resulting strong operational comparison is
\[
 T(P,f)=(G_AP,F(f)\ell_{A,P}).
\]
Indeed the displayed grade comparison and cartesian property identify
$G_A(P\otimes f^*R)$ with
$G_AP\otimes'(F(f)\ell_{A,P})^*G_BR$; the lineator compatibility
identifies the bodies. Unit and associativity are the stated coherences.
This is the established construction of \cite[equations (5.6.20)--(5.6.21)]{CM2}.

For the characterization bridge, an additional interpreted state map
$f_a:X_a\to Y_{Fa}$ must intertwine each admitted source operation $u$
with its translated operation $\widetilde u$:
$f_bu=\widetilde u f_a$. There must also be a backward translation of
\emph{admitted} labelled tests $g_a:D_{Fa}\to E_a$, satisfying
\[
 f_a(x)\models d\ \Longleftrightarrow\ x\models g_a(d),
 \qquad g_a\widetilde v_u=v_ug_b,
\]
where the continuation maps are required to be response-correct:
\[
 u(x)\models e\ \Longleftrightarrow\ x\models v_u(e),\qquad
 \widetilde u(y)\models d\ \Longleftrightarrow\ y\models\widetilde v_u(d).
\]
These acted-incidence identities make the operational lifts well-defined.
The preceding equations are the exact Chu condition with operation coherence.
It requires every pulled-back target column to be represented by an
actual source test, not merely by a semantic intersection of source columns.
Coherent label choices are stronger than columnwise representability.

\begin{proposition}[Operational and observational hypotheses together]
Under these state and test equations, $[x]\mapsto[f_a(x)]$ is a
well-defined operation-equivariant row map. Inverse image of extents has
the canonical left adjoint
\[
 (f_a)_!(A)=\cl_D(f_a[A]),\qquad (f_a)^*(B)=f_a^{-1}B.
\]
It preserves marked object concepts and intertwines the lifted operations.
The simplicial comparison has exactly the domain
\[
 K_{f_a}=\{\varnothing\}\cup
 \{\varnothing\ne U\in K_E:\exists d\in D_{Fa},\ U\subseteq C_{g_a(d)}\}.
\]
These comparisons compose, with the composite domain equal to the domain
of the first map whose image lies in the domain of the second.
\end{proposition}
\begin{proof}
Source-row equality implies equality on every $g_a(d)$, hence on every
target test after $f_a$. Preimages of target extents are intersections
of the columns $C_{g_a(d)}$, so they are source extents. This proves the
adjunction and its marked-generator property as in (CN). State and test
coherence give operation equivariance. A finite image $f_a[U]$ has a
common target witness precisely when the displayed domain condition holds.
For a second test translation $h$, the composite condition is the
existence of a witness in $\im(gh)$; this is exactly membership in the
successive domains. Closing a face preserves its common witnesses,
so the corresponding supported concept comparison has the same domain.
\end{proof}

In the display-map instance, contextual composition is an actual pullback
of spans. A functor $F$ must preserve the required display pullbacks for
a strong comparison. Otherwise the canonical map is only
\[
 F(X\times_BY)\longrightarrow FX\times_{FB}FY.
\]
With forward apex maps as span cells, this is the \emph{oplax} direction.
Neither finality of an unrelated Set-valued decomposition functor nor a
row quotient supplies its inverse. This is the exact point at which a
contextad-based application needs a new source-level proof.

The cited v2 is dated 3 November 2026 and is titled \emph{Contextful
Arrows via Wreath Products}. It is not the earlier v1 with its earlier
title and numbering. Its Section 5.4 and Remark 5.20 retain qualifications
about the free Kleisli tricocompletion and higher wreath formula. The
concrete grade calculation above does not declare those qualifications
settled, and none of the main results in this note depends on doing so.
